\begin{textsample}{POS Dim 3 – pe – Score 10.00 – pe/pe\_inf\_16.txt}  \label{ex:f3_pos_188}
2.3 CONCEPÇÕES DE \textbf{CRIANÇA}, INFÂNCIA E EDUCAÇÃO \textbf{INFANTIL} Como todas as Cantigas de Roda nascem da cultura popular, fruto da sociedade – reforçando a cultura local, a imaginação e a criatividade de um povo –, esse diálogo tem início com uma cantiga que provoca questões como : de que A Lavadeira Lava, lava lavadeira A roupinha de passear Lava, lava lavadeira A roupinha de passear Era uma menina de um tamanho assim ( 2x ) Uma trouxa de roupa assim ( 2x ) Um pedacinho de sabão assim ( 2x ) E o sol por ali assim ( 2x ) Uma lagoa desse tamanho E um pouquinho de água assim.
Uma lagoa desse tamanho E um pouquinho de água assim.
Uma lagoa desse tamanho E um pouquinho de água assim. ( Domínio Público ) \textbf{CRIANÇA} e de que INFÂNCIA este documento fala?
Em que contexto elas estão inseridas na EDUCAÇÃO \textbf{INFANTIL}?
Refletir sobre os paradoxos que envolvem as concepções de \textbf{CRIANÇA}, enquanto sujeito histórico de direitos, que se desenvolve nas interações e nas \textbf{rotinas} cotidianas a ela disponibilizadas e por ela estabelecidas com \textbf{adultos} e \textbf{crianças} nas Instituições é um grande desafio.
A estes desafios acrescentam -se ainda outras inquietações : Como a \textbf{criança} é percebida?
Qual o papel social da infância?
Qual é o significado de ser \textbf{criança} nas diferentes culturas?
Como trabalhar com a \textbf{criança} considerando seus contextos sociais, seus desenvolvimentos e formas de aprender?
Como assegurar que a educação cumpra seu papel social diante das diferenças e contradições das populações \textbf{infantis}?
Conhecer e fazer parte do universo da \textbf{criança}, considerando as diferentes idades, grupos e contextos socioculturais aos quais elas estão inseridas, desenvolver a percepção e \textbf{escuta} sensível para todas e cada uma das \textbf{crianças} é a forma mais eficaz para o professor compreender o potencial da \textbf{criança}, seus saberes e melhor forma de aprender, promovendo um planejamento real que atenda de forma prazerosa e significativa o seu desenvolvimento e aprendizagem.
Os pesquisadores como Philippe Ariès, Bernard Charlot e Walter Benjamin, através de estudos, trouxeram muitas contribuições com a inserção concreta das \textbf{crianças} e suas infâncias perante a forma de organização da sociedade em geral.
Através desses estudos, compreende -se que as \textbf{crianças} não se resumem a alguém que não é ou que virá a ser ; as \textbf{crianças} são competentes, produzem culturas e são nelas produzidas através das \textbf{brincadeiras} – que é o que as caracterizam e permitem seu poder de imaginação, fantasia e criação ; são seres indivisíveis e integrais que apresentam especificidades em todas as suas dimensões ; são sujeitos sociais e históricos, cidadãs e seres humanos detentores de direitos.
Esse modo de ver a \textbf{criança} permite entender essa fase tão importante da vida e a compreender o mundo da infância centralizado na \textbf{criança}.
Esse período complexo, repleto de desafios, não deve ser considerado apenas como uma etapa histórica, uma categoria social ou um assunto exclusivo da família que se estende do nascimento até, aproximadamente, doze anos, mas também uma pauta prioritária do Estado no sentido de garantir à \textbf{criança} direitos indivisíveis, complementáveis e inseparáveis que oferecerão a elas condições para a formação de uma pessoa integral.
Portanto, abranger a infância em toda sua magnitude exige perceber nas \textbf{crianças} a sua singularidade, o coletivo diverso do qual elas fazem parte e imergir nas diferentes culturas e saberes que produzem.
É necessário respeitar suas formas de se relacionar com o mundo e entender como se desenvolvem e aprendem, sem que o \textbf{adulto} determine o nível de desenvolvimento e aprendizagem das \textbf{crianças}.
Assim, faz -se necessário pensar como as especificidades dessa fase de vida estão sendo vivenciadas na Educação \textbf{Infantil}, na relação de \textbf{cuidado} e educação que vem se estabelecendo entre professores e \textbf{crianças}.
É preciso estabelecer uma relação de alteridade, de convivência, de interações que não sejam “ interações burocráticas ”, superficiais, pueris, e sim interações reais e dialógicas.
Interações dessa natureza ocorrem quando se estabelece uma relação respeitosa de ampla \textbf{escuta}.
Preconizadas nas DCNEI, as Interações e \textbf{Brincadeiras} correspondem aos dois grandes eixos que devem permear o currículo da Educação \textbf{Infantil}.
Vivenciar um currículo estruturado por esses dois eixos significa pensar na \textbf{criança} como um sujeito integral que se relaciona com o mundo e aprende através da mediação com o outro, com a \textbf{brincadeira} e a cultura.
Significa, ainda, compreender que a \textbf{criança} precisa estar no centro de todas as ações planejadas e promovidas para ela.
É preciso refletir até que ponto o currículo e as práticas vivenciadas pela \textbf{criança} nas instituições de Educação \textbf{Infantil} estão considerando os seus saberes e suas necessidades?
Até que ponto o protagonismo do professor com práticas adultocêntricas está considerando o olhar e a \textbf{escuta} sensível para o que a \textbf{criança} precisa experimentar.
Percebe -se que, na relação adulto/criança, para o professor o “ desafio é o de ouvi -las no que têm para nos dizer e o de escutá -las, isto é tornar as suas falas centro da compreensão dos contextos educativos e da sua transformação ” ( OLIVEIRA-FORMOSINHO ; KISHIMOTO ; PINAZZA, 2007 ).
Quando o professor permite a \textbf{criança} SER, é possível perceber que ela não se limita a aprender apenas o que ele ensina, é comum perceber que normalmente a \textbf{criança} não cede aos modelos de práticas que não a enxergam como um ser ativo e competente.
Impossível dicotomizar o \textbf{brincar} e o aprender, elas \textbf{brincam} e \textbf{brincando} aprendem o tempo todo.
Vale ressaltar que a \textbf{brincadeira}, que se fala, trata -se de uma \textbf{brincadeira} que estrutura a condição \textbf{infantil}, que oferte subsídios para a construção de \textbf{afetos} e vivências.
Ao proporcionar experiências significativas de aprendizagens e desenvolvimento para as \textbf{crianças}, o professor, incentiva sua participação num universo de \textbf{curiosidades}, descobertas, exploração de mundo e múltiplas linguagens.

% matched lemmas: adulto, afeto, brincadeira, brincar, criança, cuidado, curiosidade, escuta, infantil, rotina
\end{textsample}
